\section{Bounded co-evolution as a corollary}
\label{sec:bounded_coev}

The bounded co-evolution assumption of
\citep[Condition~(C7)]{lasser2026paper10} (invoked as condition (C7)
of the deployment-safety theorem) asserts that the per-channel
coupling magnitude $\bar{M}$ is bounded uniformly by a constant
depending on deployment-class parameters but independent of $|P|$.
This section derives the assertion from (C5.HOEFF), a new upper
exposure-rate cap (C7.RATE), and the verification protocol's
channel-projection structure under (C5.MULT).
Audit~4 of \citep[\S 8]{lasser2026paper10} previously specified an
empirical procedure for measuring $\bar{M}$ post-hoc; under the
derivation below, the audit reduces to certifying the three
structural inputs and computing $\bar{M}$ from primitives.

\subsection{Per-channel coupling magnitude (definition)}

\subsubsection{$j$-supported perturbations}

A \emph{$j$-supported signed measure} is a function
$\Delta: A \to \mathbb{R}$ with $\Delta(a) = 0$ for $a \notin A_j$
and $\sum_{a \in A_j} \Delta(a) = 0$. The set of admissible
$j$-perturbations
$\mathcal{D}_j$ is the convex subset of $j$-supported signed
measures satisfying $q_0 + \Delta \geq 0$ pointwise (so that
$q_0 + \Delta$ is a valid probability measure on $A$).

The ambient \emph{linear} space of $j$-supported zero-sum signed
measures is
\begin{equation}
\label{eq:Hj-def-companion}
  H_j \;:=\; \{\Delta : A \to \mathbb{R} \mid \Delta(a) = 0
  \text{ for } a \notin A_j, \;\; \textstyle\sum_a \Delta(a) =
  0\}.
\end{equation}
The functional $\Delta \mapsto \Delta\mu_{j'}[\Delta] := \sum_{a
\in A_{j,j'}} \Delta(a) \ell^{(j')}(a)$ extends linearly from
$\mathcal{D}_j$ to all of $H_j$.

\subsubsection{Per-step coupling magnitude}

\begin{definition}[Per-step coupling magnitude]
\label{def:M-step}
For $j \neq j'$, the \emph{per-step pairwise coupling magnitude}
of channel $j$ on channel $j'$ is the operator norm of the linear
functional $\Delta \mapsto \Delta\mu_{j'}[\Delta]$ on the linear
space $H_j$ under the $\ell_1$ norm:
\begin{equation}
\label{eq:M-step-def}
  M_{j \to j'}^{\mathrm{step}}(\mathcal{D})
  \;:=\; \sup_{\Delta \in H_j,\; \Delta \neq 0}
  \;\frac{\big| \Delta\mu_{j'}[\Delta] \big|}{\|\Delta\|_1}.
\end{equation}

The \emph{per-step deployment coupling magnitude} is
$\bar{M}^{\mathrm{step}}(\mathcal{D}) := \max_{j \neq j'}
M_{j \to j'}^{\mathrm{step}}(\mathcal{D})$.
\end{definition}

By inclusion $\mathcal{D}_j \subseteq H_j$, the operator norm on
$H_j$ upper-bounds the admissible-restricted supremum
$\sup_{\Delta \in \mathcal{D}_j \setminus \{0\}}
|\Delta\mu_{j'}[\Delta]|/\|\Delta\|_1$. We work with the
$H_j$-bound throughout; equality of the two suprema requires
$q_0$ to have strictly positive mass on every $a \in A_j$ (so
$\mathcal{D}_j$ contains a neighborhood of $0$ in $H_j$), which
we do not assume.

\subsection{Per-step closed-form bound}

\begin{lemma}[Per-step coupling bound]
\label{lem:per-step-companion}
Under (C5.HOEFF) (per-step LLR clipped to $[-\Bclip, \Bclip]$) and
the zero-outside-engagement convention, for every $j \neq j'$:
\begin{equation}
\label{eq:per-step-bound-companion}
  M_{j \to j'}^{\mathrm{step}}(\mathcal{D}) \;\leq\; \Bclip,
\end{equation}
unconditionally. When channels are action-disjoint
($A_{j, j'} = \emptyset$), $M_{j \to j'}^{\mathrm{step}} = 0$.
\end{lemma}

\begin{proof}
For any $\Delta \in H_j$ with $\Delta \neq 0$:
\begin{align*}
  \big| \Delta\mu_{j'}[\Delta] \big|
  &= \big| \sum_{a \in A_{j, j'}} \Delta(a) \ell^{(j')}(a) \big|
  \leq \sum_{a \in A_{j, j'}} |\Delta(a)| \cdot |\ell^{(j')}(a)|
  \\ &\leq \Bclip \sum_{a \in A_{j, j'}} |\Delta(a)|
  \leq \Bclip \cdot \|\Delta\|_1.
\end{align*}
Dividing by $\|\Delta\|_1 \neq 0$ and taking supremum over $H_j
\setminus \{0\}$ gives \eqref{eq:per-step-bound-companion}. When
$A_{j, j'} = \emptyset$, the sum is empty and
$\Delta\mu_{j'}[\Delta] = 0$ for all $\Delta \in H_j$, giving
$M_{j \to j'}^{\mathrm{step}} = 0$.
\end{proof}

\subsection{Cumulative cascade-window bound}

\subsubsection{Pathwise perturbation class}

A \emph{pathwise $j$-perturbation over the cascade window} is a
sequence $\{\Delta_n\}_{n=1}^{\Nev}$ where each
$\Delta_n \in H_j$, applied at SPRT exposure step $n$. Two
budget conventions are operationally relevant:
\begin{itemize}
\item \textbf{Per-step budget:} $\sup_n \|\Delta_n\|_1 \leq 1$.
\item \textbf{Total budget:} $\sum_n \|\Delta_n\|_1 \leq 1$.
\end{itemize}

The cumulative cross-channel response is
\[
  \mu_{j'}^{\mathrm{cum}}\big[\{\Delta_n\}\big]
  \;:=\; \sum_{n=1}^{\Nev(\taumeta)} \Delta\mu_{j'}[\Delta_n].
\]

\subsubsection{The new condition: upper exposure-rate cap (C7.RATE)}

\begin{assumption}[Upper exposure-rate cap, (C7.RATE)]
\label{ass:c7-rate}
There exists a deployment-class action rate cap $\lammax > 0$
independent of $|P|$, such that the SPRT exposure event count
satisfies
\begin{equation}
\label{eq:c7-rate-companion}
  \Nev(\taumeta) \;\leq\; \lammax \cdot \taumeta
\end{equation}
deterministically (operator-enforced via rate limits or action
bounds). This is distinct from (C11.CLK) of
\citep[Deployment Safety]{lasser2026paper10} which gives a lower bound on
$\Nev(\taumeta)$ for detection purposes; (C7.RATE) is an upper
bound for coupling-control purposes.

Operationally, $\lammax$ is calibrated by the deployment's rate
limits, action bounds, and channel-coupling structure (specified
in \S\ref{sec:audits}).
\end{assumption}

\subsubsection{Cumulative bound}

\begin{lemma}[Cumulative coupling bound]
\label{lem:cumulative-companion}
Under (C5.HOEFF), the zero-outside-engagement convention, and
(C7.RATE):

\emph{(a) Pairwise} cumulative coupling, per-step budget
$\sup_n \|\Delta_n\|_1 \leq 1$:
\begin{equation}
\label{eq:pairwise-cum-companion}
  \sup_{\substack{\Delta_n \in H_j \;\forall n \\
        \sup_n \|\Delta_n\|_1 \leq 1}}
  \big| \mu_{j'}^{\mathrm{cum}}\big[\{\Delta_n\}\big] \big|
  \;\leq\; \Bclip \cdot \lammax \cdot \taumeta.
\end{equation}

\emph{(b1) Aggregate-incoming} cumulative coupling under
\textbf{total per-step budget} $\sup_n \sum_j \|\Delta_n^{(j)}\|_1
\leq 1$:
\begin{equation}
\label{eq:aggregate-total-companion}
  \sup_{\substack{\Delta_n^{(j)} \in H_j \;\forall n, j \\
    \sup_n \sum_j \|\Delta_n^{(j)}\|_1 \leq 1}}
  \big| \sum_{j \neq j'}
    \mu_{j'}^{\mathrm{cum}}\big[\{\Delta_n^{(j)}\}\big] \big|
  \;\leq\; \Bclip \cdot \lammax \cdot \taumeta.
\end{equation}

\emph{(b2) Aggregate-incoming} cumulative coupling under
\textbf{per-source per-step budget} $\sup_n
\|\Delta_n^{(j)}\|_1 \leq 1$ for every $j$:
\begin{equation}
\label{eq:aggregate-source-companion}
  \sup_{\substack{\Delta_n^{(j)} \in H_j \;\forall n, j \\
    \sup_n \|\Delta_n^{(j)}\|_1 \leq 1 \;\forall j}}
  \big| \sum_{j \neq j'}
    \mu_{j'}^{\mathrm{cum}}\big[\{\Delta_n^{(j)}\}\big] \big|
  \;\leq\; (\Kch - 1) \cdot \Bclip \cdot \lammax \cdot \taumeta.
\end{equation}

The right-hand sides are deployment-class constants intensive in
$|P|$ over $\mathcal{D}$ ($\Bclip$ from (C5.HOEFF), $\lammax$
from (C7.RATE), $\taumeta$ from
\citep[Phase Redundancy]{lasser2026phase}'s scaling, $\Kch$ from
(C5.MULT)).
\end{lemma}

\begin{proof}
\emph{Part (a).} By Lemma~\ref{lem:per-step-companion},
$|\Delta\mu_{j'}[\Delta_n]| \leq \Bclip \cdot \|\Delta_n\|_1
\leq \Bclip$ for each $n$ when $\sup_n \|\Delta_n\|_1 \leq 1$.
Summing:
\[
  \big| \mu_{j'}^{\mathrm{cum}} \big| \leq
  \sum_{n=1}^{\Nev(\taumeta)} |\Delta\mu_{j'}[\Delta_n]|
  \leq \Nev(\taumeta) \cdot \Bclip
  \leq \lammax \cdot \taumeta \cdot \Bclip,
\]
the last inequality by (C7.RATE).

\emph{Part (b1) (total per-step budget).} At each step $n$,
$\sum_j \|\Delta_n^{(j)}\|_1 \leq 1$. Each per-source contribution
is bounded by $\Bclip \cdot \|\Delta_n^{(j)}\|_1$ via
Lemma~\ref{lem:per-step-companion}. Summing over $j \neq j'$ at
each step:
\[
  \big| \sum_{j \neq j'} \Delta\mu_{j'}[\Delta_n^{(j)}] \big|
  \leq \Bclip \cdot \sum_{j \neq j'} \|\Delta_n^{(j)}\|_1
  \leq \Bclip.
\]
Summing over $n$ and applying (C7.RATE) gives
\eqref{eq:aggregate-total-companion}.

\emph{Part (b2) (per-source per-step budget).} Now each source
$j$ has $\|\Delta_n^{(j)}\|_1 \leq 1$ independently. Per-source
contributions remain bounded by $\Bclip$ each; summing over
$\Kch - 1$ sources:
\[
  \big| \sum_{j \neq j'} \Delta\mu_{j'}[\Delta_n^{(j)}] \big|
  \leq (\Kch - 1) \cdot \Bclip.
\]
Summing over $n$ and applying (C7.RATE) gives
\eqref{eq:aggregate-source-companion}. Which budget convention is
operationally relevant depends on the audit specification of
admissible adversarial classes; we state both for completeness.
\end{proof}

\subsection{Bounded co-evolution as a corollary of primitive parameters}

\begin{theorem}[Bounded co-evolution corollary]
\label{thm:c7-corollary}
Under (C5.HOEFF) per-step LLR clipping, (C5.MULT) channel
multiplicity bound, (C7.RATE) upper exposure-rate cap, and the
$\taumeta$ scaling of \citep[Phase Redundancy]{lasser2026phase},
the deployment's cumulative coupling magnitude
$\bar{M}^{\mathrm{cum}}(\mathcal{D})$ over the cascade window
$\taumeta$ satisfies
\begin{equation}
\label{eq:c7-corollary}
  \bar{M}^{\mathrm{cum}}(\mathcal{D}) \;\leq\; C \cdot \Bclip
  \cdot \lammax \cdot \taumeta,
\end{equation}
where $C \in \{1, \Kch - 1\}$ depending on the budget convention
(pairwise, total, or per-source). All four factors on the RHS are
deployment-class constants intensive in $|P|$.

In particular: \citep[Condition~(C7)]{lasser2026paper10}'s
bounded-co-evolution claim that $\bar{M}$ is bounded by a
$|P|$-independent constant follows directly from
\eqref{eq:c7-corollary}, with the constant given explicitly in
terms of primitive deployment parameters.
\end{theorem}

\begin{proof}
Direct application of Lemma~\ref{lem:cumulative-companion}.
Intensivity of each RHS factor: $\Bclip$ from (C5.HOEFF) is a
deployment-class constant; $\lammax$ from (C7.RATE) is operator-
enforced and deployment-class; $\taumeta$ from
\citep{lasser2026phase} depends on deployment-class redundancy
and rate parameters but not on $|P|$; $\Kch$ from (C5.MULT) is
fixed before deployment.
\end{proof}

\subsection{Optional sharper bound under coalition-overlap stability}

The basic Theorem~\ref{thm:c7-corollary} bound uses only
$\|\Delta\|_1 \leq 1$ to get $|\Delta\mu_{j'}| \leq \Bclip$. A
sharper bound is available when the deployment imposes an
additional structural condition on the channel partition's
shared-action mass:

\begin{definition}[Optional overlap-mass stability, (C5.OVL)]
\label{def:c5-ovl}
The channel partition's shared-action mass is bounded by a
deployment-class constant $Q_{\max} < 1$, in the sense that for
every $j \neq j'$,
\[
  \sup_{\Delta \in H_j,\; \Delta \neq 0}
  \frac{\sum_{a \in A_{j, j'}} |\Delta(a)|}{\|\Delta\|_1}
  \;\leq\; Q_{\max}.
\]
\end{definition}

When (C5.OVL) is adopted, the per-step bound in
Lemma~\ref{lem:per-step-companion} sharpens to
$M_{j \to j'}^{\mathrm{step}} \leq Q_{\max} \cdot \Bclip$, and the
cumulative bound in Theorem~\ref{thm:c7-corollary} sharpens to
$\bar{M}^{\mathrm{cum}} \leq C \cdot Q_{\max} \cdot \Bclip \cdot
\lammax \cdot \taumeta$. The intensivity claim does not require
(C5.OVL); only strict-smallness ($\bar{M} \to 0$ as $Q_{\max} \to
0$, i.e., as channels become action-disjoint) requires it.

\subsection{Lemma~1 intensivity step}

\citep[Lemma~1, Step~6]{lasser2026paper10} (intensivity step in
the Goodhart-slack composition) requires
$(\epsnonresCoev)_+ \leq K_{\mathrm{coev}} \cdot \bar{M}$ to be
bounded by a $|P|$-independent constant, where $K_{\mathrm{coev}}$
is the structural co-evolution constant from
\citep{lasser2026micro}'s composition proposition.
Theorem~\ref{thm:c7-corollary} provides this directly:
\begin{equation}
  (\epsnonresCoev)_+ \;\leq\; K_{\mathrm{coev}} \cdot
  \bar{M}^{\mathrm{cum}} \;\leq\; K_{\mathrm{coev}} \cdot C \cdot
  \Bclip \cdot \lammax \cdot \taumeta,
\end{equation}
all factors deployment-class constants. The Goodhart-slack
intensivity property of \citep[Deployment Safety]{lasser2026paper10}'s deployment
claim follows.

\subsection{Failure-mode decomposition}

\citep[\S 10.2]{lasser2026paper10}'s sensitivity analysis
characterized bounded co-evolution failure as a hard-failure mode
with no graceful-degradation pathway. Under
Theorem~\ref{thm:c7-corollary}, the failure mode factors through
the failure of (C7.RATE), (C5.HOEFF), or the channel-projection
structure --- each of which has its own audit-detectable failure
signal. The hard-failure characterization is therefore replaced
by a structural diagnosis: \emph{which} primitive condition has
failed, and how the deployment can restore it.

The remainder of \citep[\S 10.2]{lasser2026paper10}'s
sensitivity analysis simplifies correspondingly: the
bounded-co-evolution paragraph becomes a forward reference to
this section, with the failure-mode taxonomy decomposed across
(C5.HOEFF), (C7.RATE), and channel-projection structure.
